% Created 2026-03-14 Sat 23:15
% Intended LaTeX compiler: pdflatex
\documentclass[11pt]{article}
\usepackage[utf8]{inputenc}
\usepackage[T1]{fontenc}
\usepackage{graphicx}
\usepackage{longtable}
\usepackage{wrapfig}
\usepackage{rotating}
\usepackage[normalem]{ulem}
\usepackage{amsmath}
\usepackage{amssymb}
\usepackage{capt-of}
\usepackage{hyperref}
\author{Prologos Research}
\date{2026-03-14}
\title{A Catalog of Lattice Structures for Propagator Networks}
\hypersetup{
 pdfauthor={Prologos Research},
 pdftitle={A Catalog of Lattice Structures for Propagator Networks},
 pdfkeywords={},
 pdfsubject={},
 pdfcreator={Emacs 30.2 (Org mode 9.7.39)}, 
 pdflang={English}}
\begin{document}

\maketitle
\setcounter{tocdepth}{3}
\tableofcontents

\section{Introduction}
\label{sec:org1d1381f}

This document catalogs lattice structures relevant to propagator networks, with
particular attention to their application in Prologos---a language unifying
dependent types, session types, linear types (QTT), logic programming, and
propagators. We organize the material into four layers:

\begin{enumerate}
\item \textbf{Fundamental lattice taxonomy}: the zoo of lattice kinds, from partial orders
through complete lattices and beyond.
\item \textbf{Concrete lattice instances}: specific lattices that arise in programming
language implementation, static analysis, and constraint solving.
\item \textbf{Composition and construction}: how lattices combine via products, sums,
lifting, Galois connections, embeddings, and functorial mappings.
\item \textbf{Applications to Prologos}: how each lattice or composition maps onto a
specific subsystem (type inference, QTT resource tracking, session types,
well-founded logic, abstract interpretation, ATMS).
\end{enumerate}

Throughout, we use standard notation: \(\sqsubseteq\) for the partial order,
\(\sqcup\) for join (least upper bound), \(\sqcap\) for meet (greatest lower bound),
\(\bot\) for bottom, \(\top\) for top.
\section{Lattice Taxonomy}
\label{sec:org43b003e}

\subsection{Partial Orders (Posets)}
\label{sec:orgd711d91}

The weakest structure. A set \(P\) with a binary relation \(\sqsubseteq\) that is
reflexive, antisymmetric, and transitive. Not every pair of elements need be
comparable (i.e., the order may be partial, not total).

\begin{itemize}
\item \textbf{Propagator relevance}: Every cell's value domain must be at least a poset.
Monotonicity of propagators is defined with respect to this order.
\end{itemize}
\subsection{Semilattices}
\label{sec:org838ea51}

\subsubsection{Join-Semilattice}
\label{sec:orgaf66c80}

A poset where every finite non-empty subset has a least upper bound (join). This
is the \textbf{minimal requirement} for a propagator cell's merge operation. The merge
laws---commutativity, associativity, idempotency---are precisely the axioms of a
join-semilattice.

\begin{itemize}
\item \textbf{Propagator relevance}: The foundational algebraic structure of the Radul-Sussman
propagator model. Every cell's merge function defines a join-semilattice.
\end{itemize}
\subsubsection{Meet-Semilattice}
\label{sec:org3c4a735}

Dual: every finite non-empty subset has a greatest lower bound (meet). Used for
descending/narrowing computations.

\begin{itemize}
\item \textbf{Propagator relevance}: Descending cells in the bilattice infrastructure
(\texttt{net-new-cell-desc}) use meet to refine from \(\top\) downward, eliminating
possibilities.
\end{itemize}
\subsection{Lattices}
\label{sec:org8469b27}

A poset that is simultaneously a join-semilattice and a meet-semilattice: every
pair of elements has both a join and a meet.
\subsubsection{Bounded Lattice}
\label{sec:org7f23fc4}

A lattice with distinguished \(\bot\) (bottom/least) and \(\top\) (top/greatest)
elements. Every finite lattice is bounded. The propagator convention: \(\bot\) = no
information, \(\top\) = contradiction.
\subsubsection{Distributive Lattice}
\label{sec:orgb00a646}

A lattice satisfying: \(a \sqcup (b \sqcap c) = (a \sqcup b) \sqcap (a \sqcup c)\)
and its dual. Distributivity is \textbf{not} automatic---the diamond lattice \(M_3\) and
the pentagon lattice \(N_5\) are the two minimal non-distributive lattices. Many
practical lattices (powersets, flat lattices, intervals over total orders) are
distributive.

\begin{itemize}
\item \textbf{Propagator relevance}: Distributivity permits certain optimizations in
constraint propagation; it guarantees that join and meet interact predictably.
\end{itemize}
\subsubsection{Modular Lattice}
\label{sec:org3f41f18}

Weaker than distributive: \(a \sqcup (b \sqcap c) = (a \sqcup b) \sqcap c\)
whenever \(a \sqsubseteq c\). Every distributive lattice is modular. The lattice of
subgroups of an abelian group is modular but not necessarily distributive.
\subsubsection{Complemented Lattice}
\label{sec:org0daf4ea}

A bounded lattice where every element \(a\) has a complement \(\bar{a}\) such that
\(a \sqcup \bar{a} = \top\) and \(a \sqcap \bar{a} = \bot\). The complement need not
be unique unless the lattice is also distributive.
\subsubsection{Boolean Algebra (Boolean Lattice)}
\label{sec:org82d50ec}

A complemented distributive lattice. Equivalently, a bounded distributive lattice
with a complement operation satisfying De Morgan's laws. The canonical example is
the powerset lattice \((\mathcal{P}(S), \subseteq)\).

\begin{itemize}
\item \textbf{Propagator relevance}: Boolean algebras model classical propositional
reasoning, ATMS worldview management, and set-based constraint domains.
\end{itemize}
\subsection{Complete Lattices}
\label{sec:orgfea4a04}

Every subset (not just finite ones) has both a join and a meet. This is the
setting for the Knaster-Tarski fixed-point theorem: every monotone function on a
complete lattice has a least and greatest fixed point.

\begin{itemize}
\item \textbf{Propagator relevance}: The Knaster-Tarski theorem is the theoretical
backbone guaranteeing that \texttt{run-to-quiescence} terminates at the least fixed
point. Completeness is needed for infinite-domain propagation (e.g., type
universes).
\end{itemize}
\subsubsection{Algebraic Lattice}
\label{sec:org06a007c}

A complete lattice where every element is the join of compact elements below it.
The lattice of ideals of a ring is algebraic. Scott domains (from domain theory)
are typically algebraic.
\subsubsection{Continuous Lattice}
\label{sec:org59d9554}

A complete lattice where every element is the directed join of elements
"way-below" it (\(x \ll y\) iff for every directed set \(D\) with \(y \sqsubseteq
\bigsqcup D\), there exists \(d \in D\) with \(x \sqsubseteq d\)). Continuous
lattices are the denotational-semantics workhorses.
\subsection{Heyting Algebras}
\label{sec:orgd961eee}

A bounded lattice with a relative pseudo-complement operation \(a \Rightarrow b\)
(the largest \(c\) such that \(a \sqcap c \sqsubseteq b\)). Every Boolean algebra is
a Heyting algebra, but not conversely. The open sets of a topological space form
a Heyting algebra under inclusion.

\begin{itemize}
\item \textbf{Propagator relevance}: Heyting algebras model intuitionistic logic, which is
the natural logic of type theory. Implications in dependent type checking
correspond to Heyting implication.
\end{itemize}
\subsection{Bilattices}
\label{sec:org9dc629b}

An algebraic structure \((B, \leq_t, \leq_k)\) with two lattice orderings:

\begin{itemize}
\item \(\leq_t\): the \textbf{truth ordering} (false \(\leq_t\) true)
\item \(\leq_k\): the \textbf{knowledge ordering} (neither \(\leq_k\) true, neither \(\leq_k\)
false, neither \(\leq_k\) both)
\end{itemize}

The four-valued Belnap bilattice \(\mathbf{FOUR}\) has elements: \(\bot_k\) (neither
true nor false / unknown), \(\mathbf{t}\) (true), \(\mathbf{f}\) (false), \(\top_k\)
(both true and false / inconsistent).

\begin{itemize}
\item \textbf{Propagator relevance}: Prologos's \texttt{bilattice.rkt} implements exactly this
structure. Bilattice variables pair an ascending cell (lower bound, knowledge
accumulates) with a descending cell (upper bound, possibilities narrow). The gap
between lower and upper yields three-valued readings: true, false, or unknown.
This is the foundation of the well-founded logic engine (WFLE).
\end{itemize}
\subsection{Quantales}
\label{sec:org30b301d}

A complete lattice equipped with an associative binary operation \(\otimes\) that
distributes over arbitrary joins: \(a \otimes (\bigsqcup S) = \bigsqcup \{a
\otimes s \mid s \in S\}\). Quantales unify lattice theory with monoid structure.

\begin{itemize}
\item \textbf{Propagator relevance}: The QTT resource semiring can be viewed through the
lens of quantales when it carries a compatible lattice ordering. Resource
usage tracking (\(0, 1, \omega\)) under the QTT semiring operations (addition
distributing over the quantity ordering) has quantale-like structure.
\end{itemize}
\section{Concrete Lattice Instances}
\label{sec:org4f95acd}

\subsection{Flat Lattices}
\label{sec:org6c7b590}

Given any set \(S\), the \textbf{flat lattice} \(S_\bot^\top\) adds a bottom \(\bot\) and top
\(\top\) with all elements of \(S\) incomparable to each other:

\begin{verbatim}
   top
 / | | \
a  b  c  d  ...
 \ | | /
   bot
\end{verbatim}

Height 2 (three levels including bot/top). Every distinct merge yields \(\top\)
(contradiction).

\begin{center}
\begin{tabular}{lll}
Instance & Elements & Application in Prologos\\
\hline
\texttt{FlatVal} & bot, val(x), top & General-purpose type inference cells\\
Constants & bot, \{c\}, top & Constant propagation\\
Symbols & bot, sym(s), top & Name resolution, identifier tracking\\
\end{tabular}
\end{center}
\subsection{Two-Point Lattice (Bool)}
\label{sec:orgc223161}

The simplest non-trivial lattice: \(\{false \sqsubseteq true\}\).

\begin{itemize}
\item Join = disjunction (OR)
\item Meet = conjunction (AND)
\item Used in: reachability analysis, boolean constraint propagation, trait
resolution flags.
\end{itemize}

Implemented as \texttt{bool-lattice} in \texttt{bilattice.rkt}.
\subsection{Three-Point Lattices}
\label{sec:org004782b}

\subsubsection{Definedness Lattice}
\label{sec:orgdf34f32}

\(\bot \sqsubseteq \mathit{defined} \sqsubseteq \top\).

Used for tracking whether a value has been computed at all, without caring about
its identity.
\subsubsection{Belnap's FOUR (Three-Point on Each Axis)}
\label{sec:orga17b6ce}

On the knowledge axis: \(\bot_k \sqsubseteq t, \bot_k \sqsubseteq f, t
\sqsubseteq \top_k, f \sqsubseteq \top_k\).

Four elements, two incomparable pairs. The bilattice of well-founded semantics.
\subsection{Powerset Lattices}
\label{sec:org710c590}

\((\mathcal{P}(S), \subseteq, \cup, \cap, \emptyset, S)\). A complete Boolean
algebra. Height = \(|S|\).

\begin{center}
\begin{tabular}{lll}
Instance & Base Set \(S\) & Application in Prologos\\
\hline
Set of constraints & Constraint IDs & Constraint store merging (union)\\
Set of assumptions & Assumption labels & ATMS environment management\\
Set of dependencies & Cell/prop IDs & Dependency tracking, worklist management\\
Set of types & Type constructors & Union type accumulation\\
Set of trait impls & Instance IDs & Trait resolution candidate sets\\
Set of capabilities & Capability tokens & Session type capability tracking\\
\end{tabular}
\end{center}
\subsection{Interval Lattices}
\label{sec:org48be46a}

\subsubsection{Numeric Intervals}
\label{sec:org35e78e3}

\([\ell, u]\) where \(\ell \leq u\). Ordered by reverse inclusion: \([a,b]
\sqsubseteq [c,d]\) iff \(c \leq a\) and \(b \leq d\) (narrower = more information).
Join = intersection.

\begin{itemize}
\item Height: infinite (for real-valued bounds)
\item Requires widening/narrowing for termination
\item Implemented in Prologos's \texttt{lattice.prologos} as the \texttt{Interval} type with
\texttt{Widenable} instance
\end{itemize}
\subsubsection{Universe Level Intervals}
\label{sec:org2d5c361}

Prologos uses universe polymorphism. Level variables form an interval lattice
where the merge narrows the range of possible levels. Zonking defaults unsolved
levels to \texttt{lzero}.
\subsection{Map/Dictionary Lattices}
\label{sec:orgd5162cf}

Given a key set \(K\) and a value lattice \(V\), the map lattice \(K \to V\) is ordered
pointwise: \(m_1 \sqsubseteq m_2\) iff \(\forall k. m_1(k) \sqsubseteq m_2(k)\).

Join and meet are computed key-by-key.

\begin{center}
\begin{tabular}{llll}
Instance & Keys & Values & Application\\
\hline
Type environment & Variable names & Type lattice & Type inference context\\
Meta solutions & Meta IDs & Term lattice & Unification state\\
QTT usage map & Variable names & Quantity lattice & Resource consumption tracking\\
Constraint store & Cell IDs & Constraint sets & Propagator network state\\
Module registry & Module names & Module contents & Namespace management\\
\end{tabular}
\end{center}

Implemented via \texttt{map-merge-with} in \texttt{lattice.prologos} and hasheq-based
registries in \texttt{constraint-cell.rkt}.
\subsection{Sign Lattice}
\label{sec:org73747f6}

\begin{verbatim}
    top
  / | \
neg zero pos
  \ | /
    bot
\end{verbatim}

Five elements. Models the sign of a numeric expression. Implemented in
\texttt{abstract-domains.prologos} as \texttt{Lattice Sign}.
\subsection{Parity Lattice}
\label{sec:orgac5deda}

\begin{verbatim}
   top
  /   \
even   odd
  \   /
   bot
\end{verbatim}

Four elements. Models even/odd parity. Implemented as \texttt{Lattice Parity} in
\texttt{abstract-domains.prologos}.
\subsection{Type Lattice}
\label{sec:org6a61c8e}

The lattice over Prologos types, where:

\begin{itemize}
\item \(\bot\) = unknown (unsolved meta)
\item \(\top\) = contradiction (inconsistent constraints)
\item Concrete types are partially ordered by the subtyping/unification relation
\item Compound types (Pi, Sigma, etc.) decompose structurally into sub-lattices
\end{itemize}

Implemented in \texttt{type-lattice.rkt}. Structural decomposition in
\texttt{elaborator-network.rkt} handles 11 compound type constructors.
\subsection{Session Type Lattice}
\label{sec:orga4deb78}

The lattice over session types, where:

\begin{itemize}
\item \texttt{sess-bot} (\(\bot\)) = unknown session
\item \texttt{sess-top} (\(\top\)) = contradictory session
\item Concrete sessions (send, recv, choice, offer, end) merge structurally
\item Continuation components merge recursively
\end{itemize}

Implemented in \texttt{session-lattice.rkt}. Behavioral subtyping (safe substitutability
of processes) induces the partial order.
\subsection{QTT Quantity Lattice}
\label{sec:org475da16}

The semiring of quantities \(\{0, 1, \omega\}\) with addition and multiplication.
Under the information ordering:

\begin{itemize}
\item \(\bot\) = unknown usage
\item \(0\) = erased at runtime
\item \(1\) = used exactly once (linear)
\item \(\omega\) = unrestricted use
\end{itemize}

The partial order: \(0 \sqsubseteq \omega\) and \(1 \sqsubseteq \omega\) (but \(0\)
and \(1\) are incomparable). This is a bounded join-semilattice.

\begin{itemize}
\item \textbf{Addition}: resource accumulation (\(0 + 1 = 1\), \(1 + 1 = \omega\))
\item \textbf{Multiplication}: scaling (\(0 \cdot x = 0\), \(\omega \cdot 1 = \omega\))
\end{itemize}
\subsection{Substitution Lattice}
\label{sec:org7861817}

The lattice of substitutions (mappings from variables to terms), ordered by
generality. The most general unifier (MGU) is the meet of all unifiers.
Robinson's unification algorithm computes this meet.

\begin{itemize}
\item \(\bot\) = the identity substitution (most general)
\item \(\top\) = failure (no unifier exists)
\item Join = anti-unification (least general generalization)
\item Meet = unification (most general unifier)
\end{itemize}
\subsection{Herbrand Interpretation Lattice}
\label{sec:org5fa162d}

For logic programming: the powerset of ground atoms, ordered by inclusion.
Immediate consequence operator \(T_P\) is monotone on this lattice, and its least
fixed point gives the minimal Herbrand model.

\begin{itemize}
\item \textbf{Propagator relevance}: Logic programming in Prologos (unification, resolution)
can be modeled as propagation over the Herbrand lattice.
\end{itemize}
\subsection{Constraint Store Lattice}
\label{sec:org1bebbdd}

The accumulation of constraints during type checking forms a lattice:

\begin{itemize}
\item \(\bot\) = empty constraint store
\item \(\top\) = inconsistent constraints
\item Join = union of constraints (more constraints = more information)
\end{itemize}

Prologos tracks these via \texttt{constraint-cell.rkt} with hasheq-based registries
using append/union merge strategies.
\subsection{ATMS Environment Lattice}
\label{sec:org758bf10}

Environments in an ATMS (sets of assumptions) form a lattice ordered by subset
inclusion. An environment \(E_1 \subseteq E_2\) means \(E_1\) is a weaker context.

\begin{itemize}
\item Moving \textbf{up} the lattice: adding assumptions, making worlds more specific
\item \textbf{Nogoods}: maximal inconsistent environments
\item \textbf{Labels}: minimal consistent environments supporting a datum
\end{itemize}

The ATMS worldview lattice enables simultaneous maintenance of multiple
hypothetical contexts without backtracking.
\section{Lattice Constructions and Compositions}
\label{sec:orgcce0dac}

\subsection{Product Lattice}
\label{sec:org678aad9}

Given lattices \(L_1, L_2, \ldots, L_n\), their \textbf{Cartesian product} \(L_1 \times
L_2 \times \cdots \times L_n\) is a lattice with componentwise operations:

\((a_1, \ldots, a_n) \sqcup (b_1, \ldots, b_n) = (a_1 \sqcup_1 b_1, \ldots, a_n \sqcup_n b_n)\)

Product lattices are the standard way to track multiple independent aspects of
program state simultaneously.

\begin{center}
\begin{tabular}{lll}
Product & Components & Application\\
\hline
Type \(\times\) Multiplicity & Type lattice, QTT quantities & Dependent type checking with QTT\\
Value \(\times\) Provenance & Value lattice, assumption set & Supported values / ATMS\\
Lower \(\times\) Upper & Ascending cell, descending & Bilattice variables\\
Sign \(\times\) Parity & Sign lattice, parity lattice & Combined numeric abstract domain\\
Type \(\times\) Session & Type lattice, session lattice & Typed communication channels\\
\end{tabular}
\end{center}
\subsection{Reduced Product}
\label{sec:orgd3898d4}

The \textbf{reduced product} of lattices \(L_1\) and \(L_2\) is a sublattice of \(L_1 \times
L_2\) that identifies redundant pairs via a reduction operator \(\rho\). The
reduction enforces logical relationships between the components: if knowing \(a_1
\in L_1\) constrains what \(a_2 \in L_2\) can be, the reduced product captures this.

\(\rho(a_1, a_2) = (\alpha_1(\gamma_1(a_1) \cap \gamma_2(a_2)), \alpha_2(\gamma_1(a_1) \cap \gamma_2(a_2)))\)

\begin{itemize}
\item \textbf{Propagator relevance}: Cross-domain propagators (\texttt{net-add-cross-domain-propagator})
implement exactly this: when information in one domain constrains another,
propagators shuttle the refinement between cells. The reduced product is the
implicit lattice of the combined network.
\end{itemize}
\subsection{Sum (Coproduct) Lattices}
\label{sec:org9f87732}

\subsubsection{Disjoint Sum}
\label{sec:orge631c8a}

\(L_1 + L_2\): tag elements with their origin, add a shared \(\bot\) and \(\top\).
Elements from different components are incomparable.
\subsubsection{Coalesced Sum}
\label{sec:org41e402b}

\(L_1 \oplus L_2\): identify the \(\bot\) elements of both lattices. Standard
construction in domain theory.

\begin{itemize}
\item \textbf{Propagator relevance}: Union types in Prologos involve a form of coalesced
sum, where different type alternatives share a common bottom (unknown).
\end{itemize}
\subsection{Lifting}
\label{sec:org854a8a3}

\(L_\bot\): add a fresh bottom element below all elements of \(L\). Dually,
\(L^\top\) adds a fresh top.

\begin{itemize}
\item \(L_\bot\): turns a lattice into one with an explicit "no information" state
\item \(L^\top\): turns a lattice into one with an explicit "contradiction" state
\item \(L_\bot^\top\): the standard "flat lattice" construction when \(L\) has no
pre-existing order
\end{itemize}

Every propagator cell domain is effectively a lifted-and-topped structure: the raw
domain \(D\) becomes \(D_\bot^\top\) where \(\bot\) = nothing and \(\top\) = contradiction.
\subsection{Smash Product}
\label{sec:orge82448c}

\(L_1 \otimes L_2\): the product with bottoms identified. In domain theory, this
corresponds to strict function spaces where \(\bot \otimes x = x \otimes \bot =
\bot\).

\begin{itemize}
\item \textbf{Propagator relevance}: When two cells must both have information before any
deduction proceeds (strict propagators), the effective domain is a smash product.
\end{itemize}
\subsection{Function Space Lattice}
\label{sec:org405e63f}

\([L_1 \to L_2]\): the set of monotone functions from \(L_1\) to \(L_2\), ordered
pointwise. If \(L_2\) is a complete lattice, so is \([L_1 \to L_2]\).

\begin{itemize}
\item \textbf{Propagator relevance}: Propagators themselves live in a function space lattice.
A propagator is a monotone function \(f : L_1 \to L_2\). The space of all such
functions forms a lattice, and composition of propagators corresponds to
function composition in this lattice.
\end{itemize}
\subsection{Powerdomain Constructions}
\label{sec:org200615f}

Powerdomains generalize powersets to domains (lattices with appropriate
continuity). Three classical variants:

\begin{center}
\begin{tabular}{llll}
Powerdomain & Operation & Models & Propagator Application\\
\hline
Hoare (lower) & Union & May-analysis (angelic nondeterminism) & Possible types, reachable states\\
Smyth (upper) & Intersection & Must-analysis (demonic nondeterminism) & Required properties, guaranteed results\\
Plotkin (convex) & Both & Combined may+must & Full abstract interpretation\\
\end{tabular}
\end{center}

\begin{itemize}
\item \textbf{Propagator relevance}: The Hoare powerdomain models angelic nondeterminism
(explored branches in logic programming); the Smyth powerdomain models demonic
nondeterminism (all branches must succeed, as in universal quantification over
session type choices).
\end{itemize}
\subsection{Galois Connections}
\label{sec:orgcd58c8f}

A \textbf{Galois connection} between posets \((C, \sqsubseteq_C)\) and \((A, \sqsubseteq_A)\)
is a pair of monotone functions:

\begin{itemize}
\item \(\alpha : C \to A\) (abstraction)
\item \(\gamma : A \to C\) (concretization)
\end{itemize}

satisfying: \(\alpha(c) \sqsubseteq_A a \iff c \sqsubseteq_C \gamma(a)\).

Equivalently, \(\alpha \circ \gamma \sqsupseteq id_A\) and \(\gamma \circ \alpha
\sqsupseteq id_C\) (overconcretizing then abstracting is sound).
\subsubsection{Properties}
\label{sec:org6555695}

\begin{itemize}
\item \textbf{Soundness}: abstraction never loses critical information
\item \textbf{Composability}: \((C \xrightarrow{\alpha_1, \gamma_1} A_1 \xrightarrow{\alpha_2, \gamma_2} A_2)\) gives \((\alpha_2 \circ \alpha_1, \gamma_1 \circ \gamma_2)\)
\item \textbf{Best abstraction}: \(\alpha\) gives the most precise element in \(A\) that
overapproximates \(c\)
\end{itemize}
\subsubsection{Galois Connection Catalog for Prologos}
\label{sec:orgc647889}

\begin{center}
\begin{tabular}{lllll}
Concrete Domain & Abstract Domain & \(\alpha\) & \(\gamma\) & Application\\
\hline
Integers \(\mathbb{Z}\) & Sign lattice & sign extraction & all ints with that sign & Numeric abstract interpretation\\
Integers \(\mathbb{Z}\) & Parity lattice & parity extraction & all ints with that parity & Parity analysis\\
Integers \(\mathbb{Z}\) & Interval \([\ell, u]\) & singleton interval & integer range & Range analysis\\
Concrete types & Type shapes & erase details & all types matching shape & Speculative type checking\\
Session protocols & Session skeletons & erase payload types & all sessions matching & Protocol conformance checking\\
Full proof terms & Proof-irrelevant tags & erase proof content & all proofs of proposition & Proof erasure (QTT \$0\$-usage)\\
Ground terms & Herbrand abstractions & generalize & instantiate & Logic programming resolution\\
Concrete states & Abstract states & widened collecting semantics & concrete state sets & Model checking via propagators\\
\end{tabular}
\end{center}

Implemented in Prologos as the \texttt{GaloisConnection} trait in \texttt{lattice.prologos} and
cross-domain propagators in \texttt{propagator.rkt}.
\subsubsection{Composition of Galois Connections}
\label{sec:orgdb4607f}

\textbf{Serial composition}: successive layers of abstraction.
\(C \xrightarrow{(\alpha_1, \gamma_1)} A_1 \xrightarrow{(\alpha_2, \gamma_2)} A_2\)
yields \((C, A_2, \alpha_2 \circ \alpha_1, \gamma_1 \circ \gamma_2)\).

\textbf{Parallel composition}: independent aspects of the same concrete domain.
\((C, A_1 \times A_2, \langle \alpha_1, \alpha_2 \rangle, \gamma_1 \sqcap \gamma_2)\).

\textbf{Propagator implementation}: serial composition = chain of cross-domain
propagators; parallel composition = reduced product with cross-domain propagators
maintaining consistency.
\subsection{Lattice Embeddings}
\label{sec:org41f5ee5}

An \textbf{embedding} \(e : L_1 \hookrightarrow L_2\) is an injective lattice homomorphism
(preserves \(\sqcup\) and \(\sqcap\)). It places \(L_1\) as a sublattice of \(L_2\).

\begin{center}
\begin{tabular}{llll}
Embedding & From & Into & Application\\
\hline
Bool \(\hookrightarrow\) FOUR & 2-valued logic & 4-valued bilattice & Classical reasoning in WFLE context\\
Flat \(\hookrightarrow\) Interval & Constants & Numeric intervals & Promoting constant to range\\
Sign \(\hookrightarrow\) Interval & Sign domain & Interval domain & Refining sign with range bounds\\
Type \(\hookrightarrow\) Type \(\times\) Q & Types & Types + quantities & Adding QTT annotations\\
Local session \(\hookrightarrow\) Global & Endpoint types & Multiparty protocol & Session type projection\\
\end{tabular}
\end{center}

Embeddings enable incremental enrichment of cell domains: start with a simple
lattice and promote to a richer one as more information arrives.
\subsection{Lattice Quotients}
\label{sec:org3451f95}

Given a lattice \(L\) and a congruence relation \(\equiv\) (an equivalence compatible
with \(\sqcup\) and \(\sqcap\)), the quotient \(L / \equiv\) is a lattice.

\begin{itemize}
\item \textbf{Propagator relevance}: Alpha-equivalence on terms, eta-equivalence, and
definitional equality all define congruences on the type lattice. Zonking
(\texttt{zonk.rkt}) can be understood as computing in a quotient lattice where
solved metas are identified with their solutions.
\end{itemize}
\subsection{Lattice Fixed-Point Operators}
\label{sec:org3732d6b}

\subsubsection{Knaster-Tarski Theorem}
\label{sec:orgeab6399}

Every monotone function \(f : L \to L\) on a complete lattice \(L\) has:
\begin{itemize}
\item A \textbf{least fixed point} \(\mu f = \bigsqcap \{x \mid f(x) \sqsubseteq x\}\)
\item A \textbf{greatest fixed point} \(\nu f = \bigsqcup \{x \mid x \sqsubseteq f(x)\}\)
\end{itemize}
\subsubsection{Kleene Iteration}
\label{sec:org47d8ee9}

\(\bot, f(\bot), f^2(\bot), \ldots\) converges to \(\mu f\) when \(L\) has finite
height or \(f\) is \$\(\omega\)\$-continuous.
\subsubsection{Widening and Narrowing}
\label{sec:org8428eae}

For infinite-height lattices, the Kleene chain may not converge. The Cousot
solution:

\begin{itemize}
\item \textbf{Widening} \(\nabla\): accelerates convergence by jumping past intermediate
values. \(a \nabla b \sqsupseteq a \sqcup b\). Guarantees termination but
overshoots.
\item \textbf{Narrowing} \(\Delta\): recovers precision. \(a \sqcap b \sqsubseteq a \Delta b
  \sqsubseteq a\). Iteratively tightens the overapproximation.
\end{itemize}

Implemented via the \texttt{Widenable} trait in \texttt{lattice.prologos} and \texttt{widen-fns}
in \texttt{propagator.rkt}.
\subsubsection{Stratified Fixed Points}
\label{sec:org0fb3032}

For logic programs with negation, the standard model is computed via
\textbf{stratification}: partition predicates into strata where negation only refers to
lower strata, compute fixed points bottom-up stratum by stratum.

Prologos implements this via the stratified quiescence protocol in the propagator
migration, where constraint resolution proceeds in ordered phases.
\subsection{Lattice-Indexed Families}
\label{sec:org7af5db9}

A family of lattices \(\{L_i\}_{i \in I}\) indexed by some set \(I\). The \textbf{dependent
product} \(\prod_{i \in I} L_i\) (pointwise operations) generalizes the simple
product. The \textbf{dependent sum} \(\sum_{i \in I} L_i\) (tagged union) generalizes the
coproduct.

\begin{itemize}
\item \textbf{Propagator relevance}: In dependent type theory, the type of a later argument
may depend on the value of an earlier one. This creates dependent lattice
structures: the lattice used for a cell may depend on the resolved value of
another cell. Prologos handles this via structural decomposition, which
lazily creates sub-cells whose domains depend on the compound type being
decomposed.
\end{itemize}
\section{Combinatorial Lattice Application Map}
\label{sec:orgda389d4}

This section maps how lattice structures combine to serve each major Prologos
subsystem.
\subsection{Type Inference Engine}
\label{sec:orgaecb7b6}

\begin{center}
\begin{tabular}{lll}
Component & Lattice Structure & Composition\\
\hline
Meta-variables & Flat lattice (unresolved/solved/conflict) & Map lattice over meta IDs\\
Type terms & Structural type lattice with decomposition & Dependent product of sub-cells\\
Constraint accumulation & Powerset lattice of constraints & Union merge\\
Trait resolution & Powerset lattice of candidate instances & Intersection for disambiguation\\
Universe levels & Interval lattice \([\ell, u]\) & Product with type lattice\\
\end{tabular}
\end{center}

Cross-cutting: Galois connection from concrete types to type shapes for
speculative checking (\texttt{save-meta-state} / \texttt{restore-meta-state!}).
\subsection{QTT Resource Tracking}
\label{sec:orgc18bb35}

\begin{center}
\begin{tabular}{lll}
Component & Lattice Structure & Composition\\
\hline
Usage quantities & \(\{0, 1, \omega\}\) semiring-lattice & Map lattice over variable names\\
Multiplicity checking & Two-point (pass/fail) lattice & Product with type environment\\
Dict param handling & Quantity lattice (mw not m0) & Embedded in type context\\
\end{tabular}
\end{center}

The semiring structure (addition + multiplication) sits atop the lattice ordering,
forming a structure akin to a quantale.
\subsection{Session Types}
\label{sec:orga896fad}

\begin{center}
\begin{tabular}{lll}
Component & Lattice Structure & Composition\\
\hline
Session skeletons & Session type lattice (send/recv/choice/end) & Structural decomposition\\
Channel capabilities & Powerset lattice of capabilities & Product with session lattice\\
Protocol conformance & Two-point (conformant/non-conformant) & Galois connection from sessions\\
Multiparty projection & Product of endpoint session lattices & Embedding from local to global\\
Duality checking & Complemented lattice on session pairs & Product with session lattice\\
\end{tabular}
\end{center}
\subsection{Well-Founded Logic Engine (WFLE)}
\label{sec:org67bef09}

\begin{center}
\begin{tabular}{lll}
Component & Lattice Structure & Composition\\
\hline
Truth values & Belnap FOUR bilattice & Product (truth \(\times\) knowledge)\\
Bilattice variables & Lower \(\times\) Upper (ascending \(\times\) descending) & Reduced product with consistency propagator\\
Predicate extensions & Powerset of ground atoms & Map lattice over predicate names\\
Negation & Antitone Galois connection on bilattice & Composes with positive fixpoint\\
Stratification & Totally ordered lattice of strata & Index for lattice-indexed family\\
\end{tabular}
\end{center}
\subsection{Abstract Interpretation Framework}
\label{sec:org36c24f9}

\begin{center}
\begin{tabular}{lll}
Component & Lattice Structure & Composition\\
\hline
Sign analysis & 5-element sign lattice & Galois connection from integers\\
Parity analysis & 4-element parity lattice & Galois connection from integers\\
Interval analysis & Interval lattice with widening & Galois connection from integers\\
Combined numeric & Reduced product (sign \(\times\) interval) & Cross-domain propagators\\
Pointer analysis & Powerset lattice of abstract locations & Map lattice over variables\\
Constant propagation & Flat lattice over constant values & Map lattice over variables\\
Reachability & Two-point lattice & Powerset over program points\\
\end{tabular}
\end{center}
\subsection{ATMS and Search}
\label{sec:org38dd8be}

\begin{center}
\begin{tabular}{lll}
Component & Lattice Structure & Composition\\
\hline
Environments & Powerset lattice of assumptions & Ordered by subset inclusion\\
Supported values & Product (value \(\times\) environment set) & Value lattice \(\times\) ATMS lattice\\
Nogood detection & Boolean lattice (consistent/inconsistent) & Map over environment lattice\\
Worldview management & Lattice of environments modulo nogoods & Quotient lattice\\
Dependency tracking & Powerset lattice of justification IDs & Map over datum lattice\\
\end{tabular}
\end{center}
\subsection{Propagator Network Infrastructure}
\label{sec:org9abc32b}

\begin{center}
\begin{tabular}{lll}
Component & Lattice Structure & Composition\\
\hline
Cell values & Domain-specific (parameterized) & Indexed family over cell IDs\\
Worklist & Powerset of propagator IDs & Priority-ordered\\
Network state & Product of all cell lattices & Massive product lattice\\
Fuel counter & Descending natural numbers & Separate from value lattice\\
Contradiction flag & Two-point lattice & Join across all cells\\
Trace/observability & Sequence lattice (append-only log) & Product with network state\\
\end{tabular}
\end{center}
\section{Advanced Composition Patterns}
\label{sec:orgea10b13}

\subsection{Lattice Transformers}
\label{sec:org403861d}

By analogy with monad transformers, we can define \textbf{lattice transformers} that
systematically enrich a base lattice:

\begin{center}
\begin{tabular}{lll}
Transformer & Type & Effect\\
\hline
\texttt{Lift\_\textbackslash{}bot} & \(L \mapsto L_\bot\) & Adds explicit "no information" bottom\\
\texttt{Lift\textasciicircum{}\textbackslash{}top} & \(L \mapsto L^\top\) & Adds explicit "contradiction" top\\
\texttt{Flat} & \(S \mapsto S_\bot^\top\) & Flattens a set into a lattice\\
\texttt{Pow} & \(L \mapsto \mathcal{P}(L)\) & Powerset over lattice elements\\
\texttt{Map\_K} & \(L \mapsto (K \to L)\) & Pointwise extension over key set\\
\texttt{Interval} & \(L \mapsto [L, L]\) & Interval abstraction (requires total order)\\
\texttt{Support} & \(L \mapsto L \times \mathcal{P}(A)\) & Adds provenance/assumption tracking\\
\texttt{Bilat} & \(L \mapsto L^{asc} \times L^{desc}\) & Bilattice pairing\\
\texttt{Widen} & \(L \mapsto L + \nabla\) & Adds widening operator for convergence\\
\end{tabular}
\end{center}

These compose: \texttt{Widen(Map\_K(Interval(\textbackslash{}mathbb\{Z\})))} gives a widened map of
integer intervals keyed by variable name---an abstract domain for interval
analysis.
\subsection{Fibered Lattices}
\label{sec:org35bd6d8}

When the lattice structure at each point depends on the value at another point,
we have a \textbf{fibered} or \textbf{dependent} lattice. In Prologos:

\begin{itemize}
\item The type of a function body depends on the type of its argument (dependent
function types / Pi types)
\item The session type of a continuation depends on the choice made at a branch
(session type choice)
\item The decomposition sub-cells created depend on the head constructor of the
compound type (structural decomposition)
\end{itemize}

These are modeled by \textbf{fibrations} in category theory: a functor \(p : \mathbf{E}
\to \mathbf{B}\) where the fiber \(p^{-1}(b)\) over each object \(b\) in the base
category is a lattice.
\subsection{Lattice Comonads and Demand}
\label{sec:orgdb5dc32}

The dual of lifting (adding bottom) is \textbf{colift} or \textbf{demand}: annotating each
lattice element with how urgently it is needed. In functional-logic programming
(narrowing), demand drives which positions to narrow first.

\begin{itemize}
\item \textbf{Needed narrowing} (Antoy/Hanus) uses definitional trees to identify demanded
positions
\item This corresponds to a comonadic structure on the lattice: \texttt{extract} gives the
current value, \texttt{extend} propagates demand downward through term structure
\end{itemize}

Prologos's FL narrowing design (\texttt{FL\_NARROWING\_DESIGN.org}) connects this to
the propagator scheduler's worklist priorities.
\subsection{Lattice Adjunctions Beyond Galois Connections}
\label{sec:org7b5b242}

Galois connections are adjunctions between poset categories. More general
adjunctions (between lattice-valued functors) enable:

\begin{itemize}
\item \textbf{Free/forgetful adjunctions}: the free lattice over a poset, the forgetful
functor from lattices to posets. Relevant to generating lattice structure from
user-defined types.
\item \textbf{Kan extensions}: generalizing Galois connections to situations where the
domains are not directly comparable but related through a third structure.
\item \textbf{Profunctor optics}: lattice-valued lenses for focusing on sub-components of
compound lattice elements (connects to structural decomposition).
\end{itemize}
\section{Speculative and Creative Applications}
\label{sec:org34053ca}

\subsection{Lattice of Lattices}
\label{sec:org5be9bc8}

The class of all finite lattices, ordered by embeddability, itself has lattice-like
structure. This meta-lattice could inform:

\begin{itemize}
\item \textbf{Domain inference}: automatically selecting the right lattice for a cell based
on the constraints feeding into it
\item \textbf{Lattice migration}: promoting a cell from a simpler lattice to a richer one
as analysis demands grow (e.g., flat \(\to\) interval \(\to\) polyhedra)
\end{itemize}
\subsection{Tropical Semiring Lattice}
\label{sec:org105c32e}

The \textbf{tropical semiring} \((\mathbb{R} \cup \{+\infty\}, \min, +)\) is a semiring
where "addition" is \(\min\) and "multiplication" is \(+\). This forms a lattice
under the \(\min\) ordering.

\begin{itemize}
\item Application: shortest-path analysis, optimality propagation, cost-aware type
checking (e.g., choosing the cheapest trait instance).
\end{itemize}
\subsection{Partition Lattice}
\label{sec:orge78abef}

The lattice of all partitions of a set, ordered by refinement. Meet = finest
common coarsening, join = coarsest common refinement.

\begin{itemize}
\item Application: equivalence class management in unification, union-find viewed
as traversal of the partition lattice.
\end{itemize}
\subsection{Concept Lattice (Formal Concept Analysis)}
\label{sec:orgb22669e}

Given a set of objects \(G\), attributes \(M\), and incidence relation \(I \subseteq G
\times M\), the \textbf{concept lattice} \(\mathfrak{B}(G, M, I)\) consists of formal
concepts (maximal rectangles in the incidence matrix).

\begin{itemize}
\item Application: discovering implicit structure in trait/type relationships,
automatic classification of types by their supported operations.
\end{itemize}
\subsection{Information Flow Lattice}
\label{sec:org16c035c}

A lattice of security levels (e.g., unclassified \(\sqsubseteq\) confidential
\(\sqsubseteq\) secret \(\sqsubseteq\) top-secret) used in information flow
analysis.

\begin{itemize}
\item Application: could model capability levels in Prologos's session types,
ensuring that session interactions respect capability hierarchies. Also
relevant to any future effect-system or region-based resource tracking.
\end{itemize}
\subsection{Lattice of Regular Languages}
\label{sec:org8ca626e}

Regular languages over an alphabet \(\Sigma\), ordered by inclusion, form a
lattice (closed under union and intersection).

\begin{itemize}
\item Application: session type protocols as regular languages, where protocol
conformance = language inclusion. Merge of session constraints = intersection
of accepted protocol languages.
\end{itemize}
\subsection{Scott Information Systems}
\label{sec:org298c143}

An information system \((A, Con, \vdash)\) where \(A\) is a set of tokens, \(Con\) is
the set of consistent subsets, and \(\vdash\) is the entailment relation. The
elements of the associated Scott domain are the ideals (downward-closed, directed,
consistent subsets).

\begin{itemize}
\item Application: a general framework for defining the information content of
cells. Each cell's domain could be specified as an information system, with
the propagator network implementing the entailment relation.
\end{itemize}
\subsection{Residuated Lattices}
\label{sec:org02e3b26}

A lattice with an additional monoid operation and left/right residuals
(generalizing division). Residuated lattices are the algebraic semantics of
substructural logics (relevant, linear, etc.).

\begin{itemize}
\item Application: direct algebraic model of Prologos's linear type system. The
residuation \(a \backslash b\) (the largest \(x\) such that \(a \otimes x
  \sqsubseteq b\)) models linear implication.
\end{itemize}
\section{Summary Table: Lattice Applicability to Prologos Subsystems}
\label{sec:org19c3dfa}

\begin{center}
\begin{tabular}{lllllllll}
Lattice & Type Inf & QTT & Sessions & WFLE & AbsInt & ATMS & Logic & Narrowing\\
\hline
Flat & X &  &  &  & X &  &  & \\
Bool (2-point) & X &  &  & X &  & X & X & \\
Powerset & X &  & X & X & X & X & X & \\
Interval &  &  &  &  & X &  &  & \\
Map/Dict & X & X &  & X & X & X &  & \\
Sign &  &  &  &  & X &  &  & \\
Parity &  &  &  &  & X &  &  & \\
Bilattice (FOUR) &  &  &  & X &  &  & X & \\
Type lattice & X &  &  &  &  &  &  & \\
Session lattice &  &  & X &  &  &  &  & \\
QTT semiring &  & X &  &  &  &  &  & \\
Substitution & X &  &  &  &  &  & X & X\\
Herbrand &  &  &  &  &  &  & X & X\\
Constraint store & X &  & X & X &  &  &  & \\
ATMS environment &  &  &  &  &  & X &  & \\
Product (general) & X & X & X & X & X & X &  & \\
Reduced product &  &  &  &  & X &  &  & \\
Galois connection &  &  & X & X & X &  & X & \\
Function space & X &  &  &  & X &  &  & \\
Heyting algebra & X &  &  &  &  &  & X & \\
Residuated lattice &  & X &  &  &  &  &  & \\
Partition lattice & X &  &  &  &  &  &  & \\
Tropical semiring &  &  &  &  &  &  &  & \\
Concept lattice & X &  &  &  &  &  &  & \\
Regular language lattice &  &  & X &  &  &  &  & \\
\end{tabular}
\end{center}
\section{References}
\label{sec:org29cb085}

\begin{itemize}
\item Radul, A. (2009). \emph{Propagation Networks: A Flexible and Expressive Substrate for Computation}. PhD thesis, MIT.
\item Sussman, G. J. \& Radul, A. (2009). "The Art of the Propagator." MIT AI Memo.
\item Cousot, P. \& Cousot, R. (1977). "Abstract Interpretation: A Unified Lattice Model for Static Analysis of Programs by Construction or Approximation of Fixpoints." \emph{POPL}.
\item Cousot, P. \& Cousot, R. (2014). "A Galois Connection Calculus for Abstract Interpretation." \emph{POPL}.
\item Kuper, L. (2014). \emph{Lattice-based Data Structures for Deterministic Parallel and Distributed Programming}. PhD thesis, Indiana University.
\item Belnap, N. (1977). "A Useful Four-Valued Logic." In \emph{Modern Uses of Multiple-Valued Logic}, Reidel.
\item Fitting, M. (1991). "Bilattices and the Semantics of Logic Programming." \emph{Journal of Logic Programming}.
\item Davey, B. A. \& Priestley, H. A. (2002). \emph{Introduction to Lattices and Order}. Cambridge University Press.
\item Garg, V. K. (2015). \emph{Introduction to Lattice Theory with Computer Science Applications}. Wiley.
\item Atkey, R. (2018). "Syntax and Semantics of Quantitative Type Theory." \emph{LICS}.
\item de Kleer, J. (1986). "An Assumption-based TMS." \emph{Artificial Intelligence}.
\item Abramsky, S. \& Jung, A. (1994). "Domain Theory." In \emph{Handbook of Logic in Computer Science}, Oxford University Press.
\item Knaster, B. (1928) and Tarski, A. (1955). "A Lattice-Theoretical Fixpoint Theorem and its Applications." \emph{Pacific Journal of Mathematics}.
\end{itemize}
\end{document}
